\documentclass[10pt,a4paper]{article}
\usepackage[margin=18mm]{geometry}
\usepackage{amsmath,amssymb,booktabs,array}
\usepackage{cite}
\renewcommand{\citepunct}{], [}
\renewcommand{\citedash}{]--[}
\input{glyphtounicode}
\pdfgentounicode=1
\usepackage[hidelinks]{hyperref}
\hypersetup{pdftitle={An urban signal-control task: model and sensitivity},
 pdfauthor={Selective-discreteness urban task review},pdfsubject={Compact methodology}}
\setlength{\parindent}{0pt}
\setlength{\parskip}{3pt}
\setlength{\abovedisplayskip}{4pt}
\setlength{\belowdisplayskip}{4pt}
\setlength{\abovedisplayshortskip}{3pt}
\setlength{\belowdisplayshortskip}{3pt}
\setlength{\emergencystretch}{2em}
\widowpenalty=10000
\clubpenalty=10000
\newcommand{\pos}[1]{[#1]_+}
\newcommand{\head}[1]{\par\vspace{4pt}\noindent\textbf{#1}\par\nobreak\vspace{1pt}}
\newenvironment{points}{\begin{list}{$\bullet$}{%
 \setlength{\leftmargin}{1.4em}\setlength{\labelwidth}{0.7em}%
 \setlength{\labelsep}{0.4em}\setlength{\itemsep}{2pt}%
 \setlength{\parsep}{0pt}\setlength{\topsep}{3pt}}}{\end{list}}
\begin{document}
\begin{center}
{\large\bfseries An urban signal-control task: model and sensitivity}
\end{center}
\vspace{-5pt}

\head{1. Research questions}
The task separates three questions about selective traffic discreteness:
\begin{points}
\item \textbf{Information:} When movement totals and physical storage agree,
does retained lane order change reachable service beyond what lawful
aggregate state and history inference can explain?
\item \textbf{Decision:} Can that distinction change which effective green
allocation reduces complete source-to-sink population time?
\item \textbf{Gradient:} Does the complete sensitivity predict beneficial
service reallocations, as checked against actual objective changes and
separate representation/optimizer comparisons?
\end{points}

\head{2. Task and global assumptions}
\textbf{Concrete task.} Control one isolated, initially symmetric four-leg
junction for a specified finite demand/route bank. One constant vector
$g\geq0$, $\sum_{p=1}^4g_p=1$, allocates effective service during the
control/release interval. The four fixed protected groups are NS
through/right, NS left, EW through/right and EW left. Loading starts
from empty; a common preparation is independent of the candidate control.
A fixed common continuation then clears every released trip.

In each starting bay approach, aggregate road $U$ feeds a physical
common lane $F$ inside the local zone, then left-only bay $B$ and shared
left/through/right continuation $T$: $U\to F\to\{B,T\}$.
The two branches have equal fork-to-stop-line length in the symmetric
core. Sections own disjoint road space and storage. Permissions may vary
between declared scenario families, but remain matched within each
comparison and order pair.

\textbf{Global modeling assumptions.} These define the traffic represented,
separately from its numerical operating domain:
\begin{points}
\item \textbf{MA1: continuous mass, fixed routes.} Quantity is continuous
vehicle-equivalent mass. Route-table references, route positions,
permissions and topology remain structural. Ordered metadata retain
realized local lane sequence; a movement grid is its mass-rank view,
not physical distance or a replacement for exact order.
\item \textbf{MA2: control acts through average service.} Green fractions
change movement service and receiving priorities. Physical capacities,
travel/wave delays, storage and saturation parameters remain fixed.
No executable within-cycle waveform is represented.
\item \textbf{MA3: processing order and physical access are distinct.}
The input record orders allocator inspections. An occupied physical lane
still imposes FIFO; an available movement can use independent access,
but a movement label cannot create a path past a resident blocker.
\item \textbf{MA4: actual admission commits the lane.} Packets balance
current legal lane loads within known transfer capability. Later packets
see earlier actual admissions. Unadmitted mass has no future-lane target
and keeps priority within its movement.
\item \textbf{MA5: causal boundary-history steps.} Sending/receiving use
pre-state cumulative histories. No immediate vacancy credit or same-section
service of new inflow is allowed. Internal head calculations produce
step totals; propagation retains affine boundary interpolation.
\item \textbf{MA6: one complete population.} Exogenous releases, routes
and preparation do not change with the candidate control. Source waiting,
moving traffic and queues are counted through clearance; fixed continuation
does not erase the state response accumulated during control.
\end{points}

For each physical section $e$, take positive finite capacity $C_e$,
free-flow/wave delays $\tau^f_e,\tau^b_e$, and storage $K_e$.
The inherited triangular-diagram and timestep compatibility is
\begin{equation}
 K_e=C_e(\tau^f_e+\tau^b_e),\qquad
 0<\Delta t\leq\min_e\{\tau^f_e,\tau^b_e\}.
 \label{eq:compact-domain}
\end{equation}
No numerical lengths, demands, minimum greens or clearance cap are
selected here. They must later instantiate this domain without changing
the traffic task.

\textbf{Why this test.} Cumulative propagation and conservative route
bookkeeping supply the physical backbone~\cite{m0,m1}. The distinctive
urban question is whether reaching a bay and exposing a shared-lane head
requires retaining ordered state. An empty bay may be unusable behind
through traffic: more left service can initially do little, while through
service can remove the obstruction. Aggregate queue/storage explanations
must be allowed to challenge this claim, and forward accuracy, correct
derivatives and useful control directions require separate evidence.

\newpage
\head{3. Traffic model}
Picture released traffic waiting at its source, travelling along $U$,
and entering the common local feeder $F$. At the fork, an accessible
left turner can choose $B$ or the shared $T$ according to current loads.
Once admitted, traffic keeps its realized lane order. Stop-line heads
consume their movement's average service and compete for downstream room.
The following components define arrival readiness, entry, discharge and
the conservative state passed to the next timestep.

\textbf{3.1. Propagation and eligible arrivals.}
At $t_k=k\Delta t$, let $N^+_{e,k},N^-_{e,k}$ be cumulative entrance/exit
and $n_{e,k}=N^+_{e,k}-N^-_{e,k}$ physical stock.
Write ${\cal H}^{\pm}_e$ for linear interpolation of stored counts:
${\cal H}^{\pm}_e((k+\theta)\Delta t)
=(1-\theta)N^\pm_{e,k}+\theta N^\pm_{e,k+1}$.
The initially empty histories are zero before loading.
The repository LTM transcription~\cite{m0} gives
\begin{align}
 A_{e,k}&=\pos{{\cal H}^+_e(t_{k+1}-\tau^f_e)-N^-_{e,k}},
 &S_{e,k}&=\min(C_e\Delta t,A_{e,k}),\label{eq:compact-sending}\\
 R_{e,k}&=\min\{C_e\Delta t,\,
 \pos{K_e+{\cal H}^-_e(t_{k+1}-\tau^b_e)-N^+_{e,k}}\}.
 \label{eq:compact-receiving}
\end{align}
Here $A$ is ready resident mass; $S,R$ are transfer allowances for the
update. They are computed before new transfers, with no future query.
Their physical-section application is the urban composition used here.

\textbf{3.2. Lane admission and accessible queues.}
A cohort keeps its route reference/position, current movement and
eligibility provenance~\cite{m1}. Consider its ready mass $q$, remaining
sending $s$, and perceived receiving-lane loads $h_\ell$, initialized
from pre-state local mass. Let $u_\ell$ be residual receiving only on
legal, physically reachable paths. The ideal balancing closure is
\begin{equation}
 V=\min(q,s,\sum_\ell u_\ell),\qquad
 x_\ell=\min\{u_\ell,[\lambda-h_\ell]_+\},\quad
 \sum_\ell x_\ell=V .
 \label{eq:compact-allocation}
\end{equation}
Only actual additions $x$ update membership, $h\leftarrow h+x$, and the
remaining sending/receiving. The remainder $q-V$ stays at its old owner,
without a fixed future-lane vector. Process cohorts in original order;
a residual prevents later cohorts of its movement from passing it in
that scan, while other movements are considered. Each sending owner has
its own blocked-movement set. These admission/scan rules are defined here
from the agreed traffic policy. Lane-entry room is distinct from junction
service: zero green does not itself close accessible queue space.

An unassigned pool inspects all ready candidates $A$ and spends at most
$S$; pre-truncating candidates to the earliest $S$ mass would recreate
cross-movement blocking. Retained cohorts keep their original eligibility.
For an actual occupied lane, access additionally requires all earlier
residents to have crossed: a physical residual in $F$ prevents later
traffic reaching either branch. The graph offers no direct $U$--$B/T$
shortcut. Thus inspecting a later movement and physically passing a
blocker remain different operations. No microscopic overtaking is added.

\textbf{3.3. Elapsed service and downstream sharing.}
For a ready stop-line head of movement $m$ on lane $\ell$, its cap is
$\mu_{\ell m}$ with physical units veh-eq/s:
\begin{equation}
 g_p\geq0,\qquad \sum_{p=1}^4g_p=1,\qquad
 \mu_{\ell m}(g)=\min\{C_\ell,s_{\ell m}g_{p(m)}\}.
 \label{eq:compact-action}
\end{equation}
At each internal head state, use downstream rate
$\rho_j=R_{j,k}/\Delta t$.
Let ${\cal L}_{ij}$ contain current heads from inbound leg $i$ to output
$j$. In this topology one leg-to-output movement has one relevant group.
Cap and reuse green-weighted leg shares, then share equally within a leg:
\begin{align}
 C_{ij}&=\sum_{\ell\in{\cal L}_{ij}}c_\ell,\quad
 w_{ij}=g_{p(i,j)},\quad W_j=\min(\rho_j,\sum_iC_{ij}),\nonumber\\
 r_{ij}&=\min(C_{ij},\theta_jw_{ij}),\quad
                    \sum_i r_{ij}=W_j,\label{eq:compact-leg-sharing}\\
 r_\ell&=\min(c_\ell,\eta_{ij}),\quad
                    \sum_{\ell\in{\cal L}_{ij}}r_\ell=r_{ij}.
 \label{eq:compact-lane-sharing}
\end{align}
Zero-cap competitors are omitted; zero receiving gives zero grants.
An inbound leg is weighted once, independent of software lane count.
This is current-ready-head sharing, not compensation for earlier
discharge to equalize final step totals.

Let $\sigma$ be elapsed internal time, $q_\ell$ residual head mass and
$S_\ell^{\rm rem}$ its remaining sending quota. Rates hold until
\begin{equation}
 \delta=\min\left\{\Delta t-\sigma,\
 \min_{\ell:r_\ell>0}
       \frac{\min(q_\ell,S_\ell^{\rm rem})}{r_\ell}\right\}.
 \label{eq:compact-head-time}
\end{equation}
Advance all active lanes by $r_\ell\delta$, process tied events together,
expose new heads and recompute sharing. Stop when all rates vanish or
time ends. A zero-rate head blocks its own lane. Constant average
receiving avoids a separate pool-exhaustion rule, but head clearance
and changing competitors still create sequential dependence.

\textbf{3.4. Conservative state and boundary records.}
Let $b_{e,k},y_{e,k}$ be accepted entrance and exit. Each internal flow
leaves one owner and enters the next with the same mass and route identity:
\begin{equation}
 N^+_{e,k+1}=N^+_{e,k}+b_{e,k},\quad
 N^-_{e,k+1}=N^-_{e,k}+y_{e,k},\quad
 n_{e,k+1}=n_{e,k}+b_{e,k}-y_{e,k}.
 \label{eq:compact-update}
\end{equation}
A source first adds its release $D_{o,k}$ to waiting $Q_{o,k}$.
A sink removes at most sending and its declared allowance.
New admissions are appended after resident service selection.
Limits $b\leq R\leq K-n$ and $y\leq S\leq n$ preserve storage and mass.
No unadmitted stock is copied into a hypothetical lane queue.

Receiving children preserve their actual admission subsequence.
For positive total entrance $b$ and entrance-rank interval $[r,r+q']$,
the affine reconstruction adapted from~\cite{m1} is
\begin{equation}
 \left[t_k+\Delta t\,\frac r b,\
       t_k+\Delta t\,\frac{r+q'}b\right],
 \qquad\text{eligibility shifted by }\tau^f_e .
 \label{eq:compact-entry}
\end{equation}
No segment is created at zero inflow. Internal section crossings keep
global route position; the actual junction crossing advances it.
Exact-compatible adjacent entries may merge without altering behavior.
The movement grid covers the full resident rank range.
This clock retains boundary histories, not new propagation-event knots.

\head{4. How the task tests information and decisions}
Compare aggregate-matched local states with different realized order.
For example, ready $F$ order $[T(0.4),L(0.6)]$, $S_F=1$ and branch
receiving $R_T=0.2$, $R_B=1$ admit only $0.2$ through: its residual
blocks left access. With the same totals, branch loads $(n_B,n_T)=(0,1.8)$
and reversed order, left uses $0.6$ bay room before through uses $0.2$.
The resulting $0.2$ versus $0.8$ total entry is a physical access contrast,
not yet a claim about different complete TST or inevitable macro failure.

With a common horizon $H$ beyond full clearance of the compared action
neighbourhood, minimize the inherited complete-time account~\cite{m0}:
\begin{equation}
 J_{\Delta t}
 =\Delta t\sum_{k=0}^{H-1}
 \left[\sum_o(Q_{o,k}+D_{o,k})+\sum_en_{e,k}\right].
 \label{eq:compact-tst}
\end{equation}
$D_{o,k}$ is available before that update's admission.
This counts source waiting and all road mass once; it is a declared
timestep quadrature, not an exact integral of every substep trajectory.
Nonclearing demand cannot disappear into a completed-trip-only score.

\begin{center}
\setlength{\tabcolsep}{5pt}
\begin{tabular}{@{}p{0.07\linewidth}p{0.86\linewidth}@{}}
\toprule
Case & Contrast held separate in the accepted suite\\
\midrule
V0 & Movement-exclusive paths: neutral matched propagation/service.\\
V1 & Shared lane: head-order effects without a finite bay.\\
V2 & Finite bay: storage with order-sensitive obstruction suppressed.\\
V3 & Same bay geometry/permissions: order and access interact.\\
V4 & Held-out variation of the admitted families and conditions.\\
\bottomrule
\end{tabular}
\end{center}
V0 neutrality requires matched service and nonbinding/equivalent
receiving; native propagation differences must be qualified rather than
credited to order. V2 must establish a sound storage explanation before
a V3 advantage supports the intended information claim.

The predictors share the physical task, action and population accounting.
B0 is aggregate; B1 uses structured movement queues; B2 adds responsive
lane/storage aggregates; B3 retains the selective ordered state.
These roles are not a universal nested information ladder.
Each macro model may receive the original fine time series through its
own converter or estimator. Declare its retained variables, memory,
observation schedule and calibration. Native subsets and successful
macro inference are legitimate; an internally retained exact ordered
ledger remains meso information. Equal current totals alone do not
establish equality of each model's full lawful history.
Separate local replay, observation-updated prediction and forecasting,
stating which observations are fixed or regenerated under an intervention.

The working native B1 source is the link queue model~\cite{wei};
B2's working direction combines lane-exit allocation~\cite{gong}
with a sourced finite-storage/access adaptation, for which~\cite{liu}
is a candidate. These native implementations still need qualification
under the agreed soundness/task-fit policy. No weaker unmatched mixture
closure is substituted for them.

First compare the same prescribed feasible service perturbations.
Then separate representation comparisons from matched-budget optimizer
comparisons. A richer evaluator uses one common fixed mapping from
effective service to executable signals, matched routes/releases and
seeds, fixed calibration and declared clearance switching.
Its full master-population time includes insertion delay plus travel
duration~\cite{sumo}; unfinished members remain accounted.
Internal improvement or finite gradients alone do not establish usefulness.

\newpage
\head{5. Complete sensitivity}
Differentiate at fixed physical parameters and structural route/order
skeleton along a feasible $g+\epsilon d$, $\sum_p d_p=0$.
A dot denotes the derivative at $\epsilon=0$.
Service changes current discharge; the resulting queues and histories
change later admission, access, spillback, source waiting and sink exit.
At a fixed state, the current TST charge and fork admission have no
direct green partial; their future response is not zero.

For an allocation with free lanes ${\cal A}$ and upper-capped lanes
${\cal U}$, differentiating the current water fill gives
\begin{align}
 \lambda&=\frac{V-\sum_{\cal U}u_\ell+\sum_{\cal A}h_\ell}
                  {|{\cal A}|},\qquad
 \dot\lambda=\frac{\dot V-\sum_{\cal U}\dot u_\ell+
                                  \sum_{\cal A}\dot h_\ell}
                  {|{\cal A}|},\nonumber\\
 \dot x_\ell&=
 \begin{cases}
 \dot\lambda-\dot h_\ell,&\ell\in{\cal A},\\
 \dot u_\ell,&\ell\in{\cal U},\\
 0,&\text{inactive lower branch}.
 \end{cases}
 \label{eq:compact-allocation-derivative}
\end{align}
Differentiate $V$ through its active minimum. Carry the result into
$h$, residual sending/receiving and $\dot q^{\rm rem}=\dot q-\dot V$
before processing the next cohort. There is no frozen-target derivative.
Fully capped branches use cap responses; stable zero branches give zero
admission without deleting residual mass sensitivity.

Green changes both service caps and current leg weights:
$\dot\mu_{\ell m}=s_{\ell m}d_{p(m)}$ on a service-limited branch
and zero on a physical-cap branch; $\dot\rho_j=\dot R_{j,k}/\Delta t$.
For uncapped leg set ${\cal F}$ and capped set ${\cal C}$,
\begin{equation}
 \dot\theta=
 \frac{\dot W-\sum_{\cal C}\dot C_{ij}
               -\theta\sum_{\cal F}\dot w_{ij}}
      {\sum_{\cal F}w_{ij}},\qquad
 \dot r_{ij}=
 \begin{cases}
 \dot C_{ij},&i\in{\cal C},\\
 w_{ij}\dot\theta+\theta\dot w_{ij},&i\in{\cal F}.
 \end{cases}
 \label{eq:compact-sharing-derivative}
\end{equation}
Within-leg sharing has the corresponding unweighted response.
With no free competitors, differentiate cap branches directly.
For an interior head interval $\delta=q/r$,
$\dot\delta=\dot q/r-q\dot r/r^2$, and
\begin{equation}
 y=\sum_a r_a\delta_a,\qquad
 \dot y=\sum_a(\dot r_a\delta_a+r_a\dot\delta_a).
 \label{eq:compact-service-derivative}
\end{equation}
The final interval's response is minus preceding duration responses.
These internal times determine timestep discharge; they are not new
propagation-history knots.

For a ready left block $q_L$ before enough through traffic and
nonbinding receiving, the selected elapsed law reduces to
\begin{equation}
 y_T=s_Tg_T\left(\Delta t-\frac{q_L}{s_Lg_L}\right),
 \qquad 0<\frac{q_L}{s_Lg_L}<\Delta t.
 \label{eq:compact-service}
\end{equation}
Unit timestep/saturations, $q_L=0.2$ and $(g_L,g_T)=(0.4,0.6)$ give
through discharge $0.3$. Moving a small service fraction from through
to left has through-discharge derivative $0.25$ on this branch.
The earlier head clears sooner, outweighing the lower through rate.
This is a local response, not a complete-time benefit claim.

Fixed delayed queries retain the response of their stored histories.
Differentiate eligibility cuts and the receiving reconstruction as well;
on a positive-inflow branch,
\begin{equation}
 \dot t_{\rm front}
 =\Delta t\left(\frac{\dot r}{b}-\frac{r\dot b}{b^2}\right),
 \label{eq:compact-entry-derivative}
\end{equation}
The tail has the analogous response. If $f$ is the fraction of a cohort's
affine eligibility segment lying before the fixed query, its ready mass
$fq$ has response $f\dot q+q\dot f$. Zero inflow creates no segment;
no epsilon repairs its denominator.

Collect histories, physical stocks, sources and ordered masses in $X_k$,
write $P_k=\dot X_k$, and let $u_k$ be the allocation applied in update $k$.
The full step and its composed response are
\begin{equation}
 X_{k+1}=F_k(X_k,u_k),\qquad
 P_{k+1}=F_{k,X}P_k+F_{k,u}\dot u_k .
 \label{eq:compact-chain}
\end{equation}
Preparation is candidate-independent. During control $\dot u_k=d$;
during fixed continuation $\dot u_k=0$, while $P_k$ keeps evolving.
No per-vehicle trajectory-gradient construction is required.
For fixed releases, total population equals cumulative release minus
sink exit $Z_k$, giving the complete accounting and derivative
\begin{equation}
 J_{\Delta t}=\Delta t\sum_{k=0}^{H-1}(D_{\leq k}-Z_k),
 \qquad
 \dot J_{\Delta t}=-\Delta t\sum_{k=0}^{H-1}\dot Z_k .
 \label{eq:compact-objective-derivative}
\end{equation}
The common horizon must clear the local action neighbourhood.
Physical access, capacity and head switches can create kinks; the formulas
apply on unchanged active signatures, with directional/one-sided analysis
at boundaries. No event-resolved history refinement or smoothing is assumed.

Sixteen exact hand-definition checks and feasibility/unique-selection
reasoning support this written model. Implementation, numerical derivative
correctness, native-rival qualification and external usefulness remain
separate evidence. The adjacent
\href{urban_methodology_specification.pdf}{detailed specification}
and \href{urban_methodology_specification.tex}{retained source}
supply complete branch rules and the source/adaptation record.

\makeatletter
\renewcommand{\section}{\@ifstar{\head}{\head}}
\makeatother
\begin{thebibliography}{6}
\setlength{\itemsep}{0pt}\setlength{\parskip}{0pt}
\bibitem{m0} Approved repository M0, \emph{LTM baseline}, Component 0.2 and
total-time account:
\href{../../../../docs/milestones/milestone_0_ltm_baseline.md}{milestone plan}.
Its recorded original basis is Yperman (2007) and UNsim v0.12.0;
the equations here use the verified repository transcription.
\bibitem{m1} Approved repository M1, \emph{Ordered route-mass ledger},
route conservation, eligibility and affine reconstruction:
\href{../../../../docs/milestones/milestone_1_ordered_route_mass_ledger.md}
{milestone plan}. Urban pool selection/access is the explicit adaptation
defined in this report.
\bibitem{wei} Wei et al., link queue model,
\href{https://arxiv.org/html/2310.12249v2}{arXiv:2310.12249v2};
journal version doi:
\href{https://doi.org/10.1016/j.trc.2024.104937}{10.1016/j.trc.2024.104937}.
\bibitem{gong} Gong et al., lane-based macroscopic node model (2025),
doi: \href{https://doi.org/10.1016/j.trb.2025.103208}
{10.1016/j.trb.2025.103208}.
\bibitem{liu} Liu and Chang, lane-group storage/blockage model,
\emph{Transportation Research C}, vol. 19, pp. 130--144, 2011,
\href{https://attap.umd.edu/wp-content/uploads/2022/05/TRC522.pdf}
{author-hosted article}; doi: 10.1016/j.trc.2010.04.005.
\bibitem{sumo} SUMO official documentation,
\href{https://sumo.dlr.de/docs/Simulation/Traffic_Lights.html}
{traffic-light programs} and
\href{https://sumo.dlr.de/docs/Simulation/Output/TripInfo.html}
{TripInfo population-time fields}.
\end{thebibliography}
\end{document}
